% !TEX program = xelatex
% 주차별 발표 자료 템플릿
% 컴파일: xelatex weekly_presentation.tex
\documentclass[aspectratio=169,11pt]{beamer}

\usepackage{fontspec}
\usepackage{kotex}
\usepackage{graphicx}
\usepackage{tikz}
\usetikzlibrary{positioning}

% Windows에 기본으로 있는 글꼴을 사용한다.
\setmainfont{Arial}
\setmainhangulfont{Malgun Gothic}
\setsansfont{Arial}
\setsanshangulfont{Malgun Gothic}

\definecolor{LectureBlue}{HTML}{0B4D99}
\definecolor{TextGray}{HTML}{303030}

\setbeamertemplate{navigation symbols}{}
\setbeamercolor{normal text}{fg=TextGray,bg=white}
\setbeamercolor{frametitle}{fg=LectureBlue,bg=white}
\setbeamercolor{title}{fg=white}
\setbeamercolor{structure}{fg=LectureBlue}
\setbeamertemplate{itemize item}{\color{LectureBlue}\small$\bullet$}
\setbeamertemplate{itemize subitem}{\color{LectureBlue}\scriptsize$\bullet$}
\setbeamertemplate{frametitle}{%
  \vspace{0.55cm}\hspace{0.48cm}\insertframetitle\par\vspace{0.12cm}}
\setbeamerfont{frametitle}{series=\bfseries,size=\LARGE}

% 기준 자료처럼 하단에 소속과 페이지 번호를 표시한다.
\setbeamertemplate{footline}{%
  \begin{tikzpicture}[remember picture,overlay]
    \fill[LectureBlue] (0,0) rectangle (0.18,0.30);
    \foreach \x in {0.22,0.31,...,0.84} {
      \draw[LectureBlue,line width=0.7pt] (\x,0) -- (\x,0.30);
    }
    \node[anchor=west,text=LectureBlue,font=\scriptsize\bfseries]
      at (0.95,0.18) {과목명 | 발표자};
    \node[anchor=east,text=gray,font=\scriptsize]
      at (15.65,0.18) {\insertframenumber};
  \end{tikzpicture}%
  \vspace{0.38cm}}

\newcommand{\week}{1주차}
\newcommand{\course}{과목명}
\newcommand{\presenter}{발표자 이름}
\newcommand{\presentationdate}{2026. 08. 26.}

\title{\course\\\week. 발표 주제}
\author{\presenter}
\date{\presentationdate}

\begin{document}

% 표지
{\setbeamertemplate{footline}{}
\begin{frame}[plain]
  \begin{tikzpicture}[remember picture,overlay]
    \fill[LectureBlue] (current page.south west) rectangle (current page.north east);
    \node[anchor=west,text=white,font=\bfseries\fontsize{27}{33}\selectfont,align=left]
      at ([xshift=1.15cm,yshift=1.3cm]current page.west) {\course\\\week. 발표 주제};
    \node[anchor=west,text=white,font=\scriptsize\bfseries]
      at ([xshift=1.15cm,yshift=-2.25cm]current page.west) {과목명 / 소속};
    \node[anchor=south east,text=white,font=\small\bfseries,align=right]
      at ([xshift=-1.1cm,yshift=1.1cm]current page.south east)
      {\presentationdate\\\presenter};
  \end{tikzpicture}
\end{frame}}

% 목차
\begin{frame}{Chapter \week}
  \begin{enumerate}
    \item 첫 번째 소주제
    \item 두 번째 소주제
    \item 세 번째 소주제
    \item 정리
  \end{enumerate}
\end{frame}

\begin{frame}{1. 첫 번째 소주제}
  \begin{itemize}
    \item 이 슬라이드에는 핵심 개념을 한 문장으로 먼저 적는다.
    \item 그 아래에 근거, 특징, 또는 예시를 2~4개 정도로 정리한다.
    \item 한 항목이 길어지면 문장을 줄이고 슬라이드를 나눈다.
  \end{itemize}
\end{frame}

\begin{frame}{핵심 예시}
  \begin{columns}[T]
    \begin{column}{0.52\textwidth}
      \begin{itemize}
        \item 왼쪽에는 설명을 둔다.
        \item 오른쪽에는 그림, 식, 파형, 또는 표를 둔다.
        \item 그림 파일은 \texttt{images/} 폴더에 넣어 사용한다.
      \end{itemize}
    \end{column}
    \begin{column}{0.42\textwidth}
      \centering
      \fbox{\rule{0pt}{4.0cm}\rule{5.2cm}{0pt}}\\[0.2cm]
      {\small 그림 또는 표 영역}
    \end{column}
  \end{columns}
\end{frame}

\begin{frame}{\week 정리}
  \begin{itemize}
    \item 핵심 내용 1
    \item 핵심 내용 2
    \item 다음 주에 이어서 볼 내용
  \end{itemize}
\end{frame}

\end{document}
